\documentclass[a4paper]{article}     
\usepackage{amsfonts}
\usepackage{amsmath}
\usepackage{amssymb}
\usepackage{graphicx}

\begin{document}

\noindent \large Auteur: Jonas Saelens \\
\noindent \large Studierichting: Informatica\\
\noindent \large Oefeningenreeks 2E nr. 11.2\\

\medskip

\normalsize

$\left(6x^2-7x+1\right)^{1-x} = \left(6x^2-7x+1\right)^x$\\

grondtal is hetzelfde:\\

$6x^2-7x+1 = 6x^2-7x+1$\\

3 gevallen: 0 of 1 of niet 0 of 1\\

geval als het 1 is:\\

$6x^2-7x+1 = 1$\\

$6x^2-7x = 0$\\

$\Delta = \left(-7\right)^2 - 4 \cdot 6 \cdot 0$\\

$\Delta = 49$\\

$x1,x2 = \dfrac{7 \pm \sqrt{49}}{12}$\\

$x1 = \dfrac{14}{12} = \dfrac{7}{6}$\\

$x2 = 0$\\

geval als het 0 is:\\

$6x^2-7x+1 = 0$\\

$\Delta = \left(-7\right)^2 - 4 \cdot 6 \cdot 1$\\

$\Delta = 49-24$\\

$\Delta = 25$\\

$x1,x2 = \dfrac{7 \pm \sqrt{25}}{12}$\\

$x1 = \dfrac{7+5}{12} = \dfrac{12}{12} = 1$\\

$x2 = \dfrac{7-5}{12} = \dfrac{2}{12} = \dfrac{1}{6}$\\

geval als het niet 0 of 1 is:\\

$1-x = x$\\

$1 = 2x$\\

$x = \dfrac{1}{2}$\\

$6\left(\dfrac{1}{2}\right)^2 - 7\left(\dfrac{1}{2}\right)+1$\\

$\Rightarrow  \dfrac{6}{4} - \dfrac{7}{2} + 1$\\

$\Rightarrow  \dfrac{3}{2} - \dfrac{7}{2} + 1$ \\

$\Rightarrow  -\dfrac{4}{2} + 1$ \\

$\Rightarrow  -\dfrac{2}{2}$ \\

$\Rightarrow  -1$ \\

$\left(-1\right)^{\tfrac{1}{2}} = \left(-1\right)^{\tfrac{1}{2}} \notin \mathbb{R}$



\begin{thebibliography}{999}
%\bibitem{a} Referentie
\end{thebibliography}
\end{document} 
